\section{Discussion}
\label{sec:discussion}

\subsection{Advice for examination boards}

Our results point to a short list of practical steps. Each is backed by the
measurements in Section~\ref{sec:results}; the numbers come from one machine
and should be read as relative, not absolute (Section~\ref{sec:threats}).

\textbf{Pre-compute every result before release.} Building each page on
request failed at twice the calibrated load; serving pre-computed results from
Redis or as static files served ten times that load with no loss, at
12.5--57.8 times less CPU per student (Section~\ref{sec:rq1}). Preparing
\num{200000} results took less than \SI{10}{s}.

\textbf{Do not count on mirrors or a cold cache.} Splitting the same hardware
into mirror sites did not help, and a cache that starts empty is worse than no
cache, because every student asks for a different result
(Sections~\ref{sec:rq1} and~\ref{sec:rq4}). If a cache is used, fill it before
release; stampede defences such as request coalescing do not help here.

\textbf{Release only when every result is written.} Students who ask before
their result exists are told it does not exist (Section~\ref{sec:rq3}).

\textbf{Update every copy when a result is corrected.} Stale results came only
from copies that were not updated: a cache without invalidation and the web
server's file cache (Section~\ref{sec:rq3}). Delete or rewrite cached copies
and keep file caches short.

\textbf{Plan for refreshes.} Students who press refresh multiply the load, and
impatient ones kept an unprotected portal down for the rest of our run
(Section~\ref{sec:refresh}).

\textbf{Use a waiting room for order, not for capacity.} A waiting room cannot
serve students faster than the backend behind it. At 5$\times$ ours did not
serve more students within the run; it turned failures into waits of a minute
or more, and in front of a fast backend it only added cost
(Section~\ref{sec:rq2}). A board that uses one should keep each student's place
across refreshes (tokens), admit in arrival order if fairness matters, and also
limit how much each admitted student may send to the backend; otherwise
impatient students overload it anyway (Sections~\ref{sec:rq2}
and~\ref{sec:refresh}). With pre-computed results, the backend took the whole
surge and no waiting room was needed.

\subsection{Threats to validity}
\label{sec:threats}

\textbf{One machine, compressed time.} All runs used one 8-core desktop machine and a
surge compressed from tens of minutes to \SI{90}{s}. We therefore report load
as multiples of a calibrated capacity and compare designs relative to each
other; absolute rates do not transfer to production hardware. The
pre-computed designs were faster than our load generator could saturate
(C1 served at least \CapCOne\ students/s, the generator's limit), so their
margins are lower bounds.

\textbf{Modelled students.} No published study measures how students refresh
on result day. Patience, refresh probabilities and delays are assumptions;
Section~\ref{sec:refresh} shows how much the conclusions depend on them.

\textbf{Modelled server cost.} The \SI{5}{ms} page-building cost stands in for
the framework overhead of real portal stacks. With \SI{2}{ms} or \SI{0}{ms},
A and B0 took the 2$\times$ load easily (Section~\ref{sec:rq1}), so absolute
capacities depend on it; the cost order of the designs did not.

\textbf{Scope.} We do not model login or captcha services, bots, denial of
service, DNS, or a real CDN, each of which has caused real result-day outages.
Results are synthetic; three repetitions per configuration give wide
intervals for some metrics.

